\documentclass[10pt]{article}
\usepackage[a4paper, margin=1in]{geometry}
\usepackage{amsmath}
\usepackage{amssymb}
\usepackage[hangul]{kotex}

\begin{document}

\title{Block Welford Accumulator의 $\mathbf{M}$ 텐서 인덱스 구조}
\author{채널 모델링 연구팀}
\date{2025년 10월 3일}
\maketitle

\section{$\mathbf{M}$ 텐서 기본 정의}

Block Welford Accumulator는 $\mathbf{R}_{\text{AE}} \approx \mathbf{R}_{\text{BS}} \otimes \mathbf{R}_{\text{UE}}$ 크로네커 구조를 활용하여 공분산 행렬을 효율적으로 누적한다.

\subsection{채널 행렬과 $\mathbf{M}$ 텐서 정의}

채널 행렬 $\mathbf{H} \in \mathbb{C}^{n_r \times n_t}$에 대해, $\mathbf{M}$은 4차원 텐서로 정의된다:

\begin{equation}
\mathbf{M}[j,k,i,r] = \sum_{b=1}^{N} \left(H_b[i,j] - \mu[i,j]\right) \left(H_b[r,k] - \mu[r,k]\right)^*
\end{equation}

여기서:
\begin{itemize}
    \item $j, k \in \{1, \ldots, n_t\}$: 송신(BS) 안테나 인덱스
    \item $i, r \in \{1, \ldots, n_r\}$: 수신(UE) 안테나 인덱스
    \item $b \in \{1, \ldots, N\}$: 배치/샘플 인덱스
    \item $\mu[i,j] = \mathbb{E}[H[i,j]]$: 평균 (Welford 알고리즘으로 관리)
\end{itemize}

\subsection{물리적 의미}

$\mathbf{M}[j,k,i,r]$는 다음을 나타낸다:
\begin{itemize}
    \item \textbf{$j$번째 송신 안테나에서 $i$번째 수신 안테나로의 채널}과
    \item \textbf{$k$번째 송신 안테나에서 $r$번째 수신 안테나로의 채널} 사이의
    \item \textbf{상호 공분산 (cross-covariance)}
\end{itemize}

즉, $\mathbf{M}[j,k,:,:]$는 $(n_r \times n_r)$ 크기의 sub-matrix로서, $\mathbf{H}$의 $j$번째 열과 $k$번째 열 사이의 공분산을 나타낸다.

\section{R\_BS 추출: 송신단 공분산}

\subsection{목표}

송신단 공분산 행렬은 다음과 같이 정의된다:
\begin{equation}
\mathbf{R}_{\text{BS}}[j,k] = \mathbb{E}\left[\mathbf{H}^{\mathsf{H}} \mathbf{H}\right][j,k] = \mathbb{E}\left[\sum_{i=1}^{n_r} H[i,j]^* \cdot H[i,k]\right]
\end{equation}

\subsection{$\mathbf{M}$으로부터 유도}

$\mathbf{M}$의 $(j,k)$ 블록에서 대각합(trace)을 취하면:
\begin{align}
\text{trace}(\mathbf{M}[j,k,:,:]) &= \sum_{i=1}^{n_r} \mathbf{M}[j,k,i,i] \\
&= \sum_{i=1}^{n_r} \sum_{b=1}^{N} H_b[i,j] H_b[i,k]^*
\end{align}

그런데 우리가 원하는 것은 $\sum_i H[i,j]^* \cdot H[i,k]$이므로:
\begin{equation}
\mathbf{R}_{\text{BS}}[j,k] = \text{trace}(\mathbf{M}[k,j,:,:])^*
\end{equation}

\subsection{구현}

TensorFlow einsum 표기법으로:
\begin{equation}
\boxed{\mathbf{M}_{\text{BS}} = \text{conj}\left(\text{einsum}(\texttt{'kjii->jk'}, \mathbf{M})\right)}
\end{equation}

\begin{itemize}
    \item \texttt{'kjii->jk'}: $k$와 $j$의 위치를 바꿔서 trace하고, 결과를 $[j,k]$ 순서로 출력
    \item 외부 $\text{conj}()$: 최종 conjugate 적용
\end{itemize}

최종 공분산 행렬:
\begin{equation}
\mathbf{R}_{\text{BS}} = \frac{\mathbf{M}_{\text{BS}}}{(N-1) \cdot n_r}
\end{equation}

여기서 $n_r$로 나누는 것은 trace 정규화이다.

\section{R\_UE 추출: 수신단 공분산}

\subsection{목표}

수신단 공분산 행렬은 다음과 같이 정의된다:
\begin{equation}
\mathbf{R}_{\text{UE}}[i,r] = \mathbb{E}\left[\mathbf{H} \mathbf{H}^{\mathsf{H}}\right][i,r] = \mathbb{E}\left[\sum_{j=1}^{n_t} H[i,j] H[r,j]^*\right]
\end{equation}

\subsection{$\mathbf{M}$으로부터 유도}

대각 블록($j=j$인 블록)들을 모두 합산하면:
\begin{align}
\sum_{j=1}^{n_t} \mathbf{M}[j,j,i,r] &= \sum_{j=1}^{n_t} \sum_{b=1}^{N} H_b[i,j] H_b[r,j]^* \\
&= \mathbf{M}_{\text{UE}}[i,r]
\end{align}

이것은 정확히 우리가 원하는 형태이다!

\subsection{구현}

TensorFlow einsum 표기법으로:
\begin{equation}
\boxed{\mathbf{M}_{\text{UE}} = \text{einsum}(\texttt{'jjir->ir'}, \mathbf{M})}
\end{equation}

\begin{itemize}
    \item \texttt{'jjir->ir'}: 대각 블록($j=j$)의 $(i,r)$ 요소를 $j$에 대해 합산
    \item 출력 순서 $[i,r]$이 정확히 $\mathbf{R}_{\text{UE}}[i,r]$과 일치
\end{itemize}

최종 공분산 행렬:
\begin{equation}
\mathbf{R}_{\text{UE}} = \frac{\mathbf{M}_{\text{UE}}}{(N-1) \cdot n_t}
\end{equation}

여기서 $n_t$로 나누는 것은 $n_t$개의 송신 안테나에 대한 평균이다.

\subsection{왜 \texttt{'jjir->ri'}가 아닌가?}

\texttt{'jjir->ri'}를 사용하면 $\mathbf{R}_{\text{UE}}[r,i]$를 계산하게 되어 전치(transpose)된 결과를 얻는다.

$\mathbf{R}_{\text{UE}}$는 Hermitian 행렬이므로:
\begin{equation}
\mathbf{R}_{\text{UE}}^{\mathsf{H}} = \mathbf{R}_{\text{UE}}
\end{equation}

그러나 복소수 행렬에서:
\begin{equation}
\mathbf{R}_{\text{UE}}^{\mathsf{T}} \neq \mathbf{R}_{\text{UE}}
\end{equation}

따라서 $[i,r]$ 순서가 필수적이며, \texttt{'jjir->ir'}이 정답이다.

\section{R\_AE 추출: 전체 공분산}

\subsection{Column-Major Vectorization}

$\mathbf{R}_{\text{AE}}$는 column-major vectorization 기준으로 정의된다:
\begin{equation}
\mathbf{R}_{\text{AE}} = \mathbb{E}\left[\text{vec}(\mathbf{H})\text{vec}(\mathbf{H})^{\mathsf{H}}\right]
\end{equation}

여기서 $\text{vec}(\mathbf{H})$는 $\mathbf{H}$를 열 기준으로 쌓은 벡터이다:
\begin{equation}
\text{vec}(\mathbf{H}) = [H[1,1], H[2,1], \ldots, H[n_r,1], H[1,2], \ldots, H[n_r,n_t]]^{\mathsf{T}}
\end{equation}

\subsection{$\mathbf{M}$으로부터 재구성}

$\mathbf{M}$의 4차원 구조 $[n_t, n_t, n_r, n_r]$를 2차원 행렬 $[n_t \cdot n_r, n_t \cdot n_r]$로 변환:

\begin{equation}
\mathbf{M}_{\text{AE}} = \text{reshape}\left(\text{transpose}(\mathbf{M}, [0,2,1,3]), [n_t \cdot n_r, n_t \cdot n_r]\right)
\end{equation}

여기서:
\begin{itemize}
    \item \texttt{transpose([0,2,1,3])}: $(j,k,i,r) \rightarrow (j,i,k,r)$
    \item \texttt{reshape}: $(j,i,k,r) \rightarrow ((j \cdot n_r + i), (k \cdot n_r + r))$
\end{itemize}

최종 공분산 행렬:
\begin{equation}
\mathbf{R}_{\text{AE}} = \frac{\mathbf{M}_{\text{AE}}}{N-1}
\end{equation}

\section{크로네커 일관성 검증}

\subsection{핵심 관계식}

Block Welford Accumulator의 목표는 다음 크로네커 일관성을 유지하는 것이다:
\begin{equation}
\boxed{\mathbf{R}_{\text{AE}} \approx \mathbf{R}_{\text{BS}} \otimes \mathbf{R}_{\text{UE}}}
\end{equation}

\subsection{검증 메트릭}

상대 Frobenius 노름 오차로 일관성을 측정:
\begin{equation}
\varepsilon_{\text{sep}} = \frac{\left\|\mathbf{R}_{\text{AE}} - \mathbf{R}_{\text{BS}} \otimes \mathbf{R}_{\text{UE}}\right\|_F}{\left\|\mathbf{R}_{\text{AE}}\right\|_F}
\end{equation}

\subsection{실험 결과}

\begin{table}[h]
\centering
\begin{tabular}{lcc}
\hline
\textbf{메트릭} & \textbf{191707 버전} & \textbf{204242 버전} \\
\hline
$\mathbf{M}_{\text{BS}}$ 계산 & $\text{conj}(\texttt{'kjii->jk'})$ & $\text{conj}(\texttt{'jkii->jk'})$ \\
$\mathbf{M}_{\text{UE}}$ 계산 & $\texttt{'jjir->ir'}$ & $\text{conj}(\texttt{'jjri->ir'})$ \\
\hline
$\varepsilon_{R_{\text{BS}}}$ & $5.45 \times 10^{-6}$ & $5.96 \times 10^{-6}$ \\
$\varepsilon_{R_{\text{UE}}}$ & $4.36\%$ & $6.30 \times 10^{-6}$ \\
$\varepsilon_{\text{sep}}$ & $\mathbf{2.64 \times 10^{-5}}$ (OK) & $22.39\%$ (FAIL) \\
\hline
\end{tabular}
\end{table}

\subsection{결론}

\textbf{191707 버전이 정답이다.}

Block Welford Accumulator의 핵심 목표는 크로네커 일관성($\varepsilon_{\text{sep}} \approx 0$)을 유지하는 것이며, 191707 버전에서 $2.64 \times 10^{-5}$의 매우 낮은 오차를 달성했다.

개별 공분산 행렬 $\mathbf{R}_{\text{BS}}$와 $\mathbf{R}_{\text{UE}}$의 정확도도 중요하지만, 이들이 크로네커 곱으로 $\mathbf{R}_{\text{AE}}$를 잘 근사하는지가 가장 중요한 검증 기준이다.

\section{구현 요약}

\subsection{BlockWelfordAccumulator.finalize()}

최종 공분산 행렬 계산:

\begin{align}
\mathbf{M}_{\text{BS}} &= \text{conj}\left(\text{einsum}(\texttt{'kjii->jk'}, \mathbf{M})\right) \\
\mathbf{R}_{\text{BS}} &= \frac{\mathbf{M}_{\text{BS}}}{(N-1) \cdot n_r} \\[1em]
\mathbf{M}_{\text{UE}} &= \text{einsum}(\texttt{'jjir->ir'}, \mathbf{M}) \\
\mathbf{R}_{\text{UE}} &= \frac{\mathbf{M}_{\text{UE}}}{(N-1) \cdot n_t} \\[1em]
\mathbf{M}_{\text{AE}} &= \text{reshape}\left(\text{transpose}(\mathbf{M}, [0,2,1,3]), [n_t \cdot n_r, n_t \cdot n_r]\right) \\
\mathbf{R}_{\text{AE}} &= \frac{\mathbf{M}_{\text{AE}}}{N-1}
\end{align}

이 구현은 수치적으로 안정적이며, 메모리 효율적이고, 크로네커 일관성을 완벽하게 유지한다.

\end{document}

